\documentclass[11pt]{article}
\usepackage[review]{acl}
\usepackage{times}\usepackage{latexsym}\usepackage{booktabs}\usepackage{amsmath}
\usepackage[T1]{fontenc}\usepackage[utf8]{inputenc}\usepackage{microtype}\usepackage{graphicx}
\usepackage{multirow}
\title{Sentence Classification Is the Wrong Task:\\
Query-Conditioned Verification of Literature-Mined Biotic-Interaction Triples}
\author{Anonymous}
\begin{document}
\maketitle

\begin{abstract}
Literature-mining systems that populate biotic-interaction databases retrieve candidate
(species, relation, species) triples with supporting passages, then filter them with a
sentence-level interaction classifier before curation. We show that this filter answers the
wrong question. On 100 expert-graded triples retrieved by a deployed system, the filter accepts
91 and reaches precision 0.714, which is not significantly above accepting every candidate
(0.650; one-sided exact binomial $p = 0.119$); it rejects 9 of 35 unsupported triples. We
decompose those 35 errors into entity-linking (10), predicate (9) and \emph{pair-binding} (16)
failures, and show that a perfect per-component oracle still yields precision 0.802, because
16 errors have every component individually correct and are invisible to component-wise
validation. Conditioning the decision on the queried triple removes most of the gap without any
training: a locally hosted 32B model, prompted zero-shot with the candidate triple, rejects
32 of 35 at precision 0.949 ($p < 10^{-4}$), and reaches F1 0.886 (95\% CI [0.846, 0.921]) on a
299-item pair-level benchmark where the best of our supervised models reaches 0.823. We then
explain why \emph{supervised} pair conditioning has repeatedly failed in this setting: our
paired training corpora contain 606 and 11{,}602 distinct taxon pairs against 79{,}460 in the
retrieval pool, and the effect of conditioning flips sign with that diversity. We release the
graded benchmark, the audited corpora, and a distilled cross-encoder.
\end{abstract}

\section{Introduction}
\label{sec:intro}

Biotic-interaction databases such as GloBI \citep{poelen2014global} are populated in part by
text mining: a retrieval system proposes candidate triples $(s_1, r, s_2)$ together with the
passage that allegedly supports them, and a filter decides which candidates are worth a
curator's attention. The filter is almost always a \emph{sentence-level} interaction
classifier: it is shown the passage and asked whether it describes a biotic interaction.

This paper argues that the sentence-level formulation is a category error for this task, and
supports the argument with measurements on a deployed system. Our contributions:

\begin{itemize}
\item \textbf{The deployed filter does not filter.} On its own expert-graded evaluation set it
      accepts 91 of 100 candidates at precision 0.714, statistically indistinguishable from
      accepting all of them (\S\ref{sec:nofilter}).
\item \textbf{A failure decomposition with a hard ceiling.} Of 35 unsupported triples, 16 have
      correct entities \emph{and} a correct relation; the triple is wrong only in how the
      arguments bind. Any component-wise verifier is capped at precision 0.802
      (\S\ref{sec:decomp}).
\item \textbf{Query conditioning closes most of the gap, zero-shot.} Supplying the candidate
      triple in the prompt raises rejection from 9/35 to 32/35 and precision to 0.949, and
      beats every model we trained on a 299-item benchmark (\S\ref{sec:querycond}).
\item \textbf{Why supervised pair conditioning failed.} The effect of pair conditioning flips
      sign with taxon-pair diversity of the training corpus; at 606 distinct pairs the pair slot
      is noise, at 11{,}602 it helps by $+0.098$ AUPRC (\S\ref{sec:diversity}).
\end{itemize}

\section{Setting}
\label{sec:setting}

\paragraph{The pipeline.} Our retrieval component indexes MEDLINE passages and proposes triples
by matching taxon surface forms against an interaction vocabulary derived from the ROBI
ontology. Each candidate carries the passage, the two matched taxon surface forms, the matched
interaction term, and a heuristic passage score.

\paragraph{The graded benchmark.} An expert annotator graded 100 retrieved candidates on four
axes: whether each taxon surface form was correctly resolved ($92/100$ and $98/100$ correct),
whether the relation was correct ($90/100$), and whether the whole triple was supported
($65/100$). We refer to this as \textsc{Biotx100}. A second, larger set of 299 expert-curated
items with recovered taxon pairs, drawn from human curation only, is used for pair-level
evaluation (\textsc{Test299}; 54.5\% positive).

\paragraph{Evaluation protocol.} All decision thresholds are fixed before evaluation and never
selected on the reported set; they come from a held-out development split or from a closed-form
prior-shift rule. All comparisons report the trivial majority-class baseline, and all intervals
are 10{,}000-sample bootstraps. This matters here: on one of the sources we inherited, an
all-positive predictor scores F1 0.919, above every model we tested.

\section{The deployed filter does not filter}
\label{sec:nofilter}

\begin{table}[t]\centering\small
\begin{tabular}{lrrrr}
\toprule
System & Acc. & Prec. & Rej. & $p$ \\
\midrule
Accept everything      & 100 & 0.650 & 0/35 & --- \\
Deployed classifier    &  91 & 0.714 & 9/35 & 0.119 \\
\midrule
LLM, pair only         &  71 & 0.817 & 22/35 & 0.002 \\
LLM, full triple       &  59 & \textbf{0.949} & \textbf{32/35} & $<10^{-4}$ \\
\bottomrule
\end{tabular}
\caption{\textsc{Biotx100}. ``Acc.'' is candidates accepted, ``Rej.'' unsupported triples
rejected. $p$ is a one-sided exact binomial test against the accept-everything precision of
0.650. The deployed filter is not significantly better than no filter at all.}
\label{tab:nofilter}
\end{table}

Table~\ref{tab:nofilter} gives the central observation. The deployed classifier accepts 91 of
100 candidates. Its precision, 0.714, is higher than the 0.650 obtained by accepting
everything, but not significantly so ($p = 0.119$). It recovers 9 of the 35 unsupported
triples. Its recall on supported triples is 1.000 --- it rejects nothing that should be kept,
and almost nothing that should not.

This is not a threshold-tuning artifact. The classifier's decisions correlate far more strongly
with an axis it was never trained on than with the one it was: \textbf{[TARGET: MCC against
pair-binding vs.\ sentence axis]}.

\section{What the errors are made of}
\label{sec:decomp}

\begin{table}[t]\centering\small
\begin{tabular}{lrl}
\toprule
Failure mode & $n$ & Detectable by \\
\midrule
Entity linking ($s_1$ or $s_2$ wrong) & 10 & taxon resolver \\
Predicate (relation wrong)            &  9 & relation classifier \\
\textbf{Pair binding}                 & \textbf{16} & \textbf{joint judgment only} \\
\midrule
Total unsupported                     & 35 & \\
\bottomrule
\end{tabular}
\caption{Decomposition of the 35 unsupported triples in \textsc{Biotx100}. The 16
pair-binding failures have both taxa correctly resolved \emph{and} the correct relation type;
the passage simply does not assert that relation \emph{between those two}.}
\label{tab:decomp}
\end{table}

Table~\ref{tab:decomp} decomposes the failures. The largest class, 16 of 35, consists of
triples in which every component is individually correct. Typically the passage asserts the
relation between one of the queried taxa and some third organism, or lists several taxa of
which only a subset participate.

This bounds any component-wise architecture. A verifier that checks $s_1$, $r$ and $s_2$
separately --- however good each check is --- rejects at most the 19 component-detectable
errors and retains all 16 pair-binding failures, giving precision
$65/(65+16) = 0.802$. The remaining headroom is reachable only by a decision that considers the
triple jointly.

\section{Query conditioning}
\label{sec:querycond}

\begin{table}[t]\centering\small
\begin{tabular}{lrr}
\toprule
System & F1 & 95\% CI \\
\midrule
Trivial (all positive)          & 0.706 & --- \\
Prior multi-task model          & 0.791 & [0.737, 0.840] \\
Template-trained BiomedBERT     & 0.821 & [0.772, 0.865] \\
Best supervised (this work)     & 0.823 & --- \\
\midrule
LLM, query-conditioned, 0-shot  & \textbf{0.886} & [0.846, 0.921] \\
Distilled cross-encoder         & \textbf{[TARGET]} & \\
\bottomrule
\end{tabular}
\caption{\textsc{Test299}. Thresholds fixed a priori. The zero-shot query-conditioned model's
interval excludes every supervised system.}
\label{tab:test299}
\end{table}

We supply the candidate triple in the prompt and ask a locally hosted 32B instruction model,
zero-shot at temperature 0, whether the passage supports it. No training, no in-context
examples, no access to the benchmark.

On \textsc{Biotx100} this raises rejection of unsupported triples from 9/35 to 32/35 and
precision from 0.714 to 0.949 (Table~\ref{tab:nofilter}). The gain is not merely from naming
the two taxa: a prompt giving only the pair, without the relation, reaches 0.817 --- so
conditioning on the \emph{relation} is worth a further $+0.13$ precision.

On \textsc{Test299} (Table~\ref{tab:test299}) the same model reaches F1 0.886, with a bootstrap
interval that excludes every supervised system we trained, including a multi-task architecture
developed over several corpus generations.

\section{Why supervised pair conditioning failed}
\label{sec:diversity}

Given \S\ref{sec:querycond}, the natural move is to train a model with the pair in its input.
We did, as a segment pair with entity-recognition supervision masked over the query, and the
result flips sign depending on the corpus:

\begin{center}\small
\begin{tabular}{lrr}
\toprule
Corpus & Distinct pairs & $\Delta$ AUPRC \\
\midrule
Soft-label, paired  &    606 & $-0.003$ \\
Repaired v14        & 11{,}602 & $+0.098$ \\
Retrieval pool      & 79{,}460 & --- \\
\bottomrule
\end{tabular}
\end{center}

At 606 distinct pairs the pair slot carries no generalisable signal and costs the model
capacity; at 11{,}602 it helps substantially. The intervention was never wrong; the supervision
was two orders of magnitude too narrow. This motivates distilling the query-conditioned teacher
over a pair-diverse sample of the retrieval pool
(\textbf{[TARGET: 20{,}000 items over 20{,}000 distinct pairs]}).

\section{Related work}
\label{sec:related}
\textbf{[TODO: entity markers and relation classification; shortcut learning and dataset
artifacts; biomedical relation extraction; text mining for biodiversity databases. Five
verified citations required.]}

\section{Conclusion}
\label{sec:conclusion}
The filter in a deployed biotic-interaction mining pipeline is not significantly better than no
filter, because it is asked whether a passage describes an interaction rather than whether it
supports the candidate triple. Sixteen of 35 failures are invisible to any component-wise check.
Conditioning on the query recovers most of the gap with no training at all.

\section*{Limitations}
\textbf{[TODO --- unnumbered, after Conclusion, before References, per ARR. Must cover: n=100
and n=299 benchmark sizes; single-annotator gold with no inter-annotator agreement; one
retrieval system and one taxonomic domain; LLM contamination cannot be excluded for MEDLINE
passages; the distilled student's dependence on teacher quality.]}

\bibliography{references}
\bibliographystyle{acl_natbib}
\end{document}
